\documentclass[11pt]{article}
\usepackage[margin=25mm]{geometry}
\usepackage{fontspec}
\setmainfont{Latin Modern Roman}
\usepackage{amsmath,amssymb,mathtools}
\usepackage{unicode-math}
\setmathfont{Latin Modern Math}
\usepackage{xurl,hyperref}
\hypersetup{colorlinks=true,urlcolor=blue}
\setcounter{secnumdepth}{0}
\setlength{\parindent}{0pt}
\setlength{\parskip}{5pt}
\setlength{\emergencystretch}{3em}
\newcommand{\sref}[2]{\hyperlink{source:#1:#2}{[S#2]}}
\newcommand{\cataloguescope}[1]{\paragraph{Recorded target.}#1}
\begin{document}
\section{Mahler's Conjecture}\label{statement}

\subsection{Definitions and mathematical statement}

\paragraph{Shared notation.}$\mathbb N=\{1,2,\ldots\}$, and $\mathbb Z,\mathbb Q,\mathbb R,\mathbb C$ denote the integers, rationals, reals and complex numbers. Any other conventions are specified locally.


A convex body \(K\subset\mathbb R^n\) is a compact convex set with
nonempty interior.  In the symmetric formulation, its center is moved
to the origin, so \(K=-K\).  Its polar body is
\[
 K^\circ=\{y\in\mathbb R^n:\langle x,y\rangle\le1
                           \text{ for every }x\in K\},
\]
where \(\langle\cdot,\cdot\rangle\) is the Euclidean inner product.
Write \(|K|\) for its \(n\)-dimensional volume.  The cube
\([-1,1]^n\) and its polar, the cross-polytope
\(\{x:\sum_i|x_i|\le1\}\), have volume product \(4^n/n!\).
\paragraph{Statement recorded in the supplied source.}
For every \(n\ge1\) and every origin-symmetric convex body
\(K\subset\mathbb R^n\), must
\[
 |K|\,|K^\circ|\ge\frac{4^n}{n!}?
\]
Equivalently, do the cube and the cross-polytope attain the smallest
possible volume product in the centrally symmetric class? \sref{statement}{2}
\subsection{Short English statement}
Does every centrally symmetric convex body have volume product at least that of a cube, equivalently that of its polar cross-polytope?
\subsection{Sources}
\begin{description}
\item[\hypertarget{source:statement:1}{S1}] ULAM.AI and the UnsolvedMath Contributors, \emph{UnsolvedMath: A Curated Collection of Open Mathematics Problems}, catalogue entry. CC BY 4.0. \url{https://huggingface.co/datasets/ulamai/UnsolvedMath}.
\item[\hypertarget{source:statement:2}{S2}] Hiroshi Iriyeh and Masataka Shibata, \emph{Symmetric Mahler's conjecture for the volume product in the three dimensional case}, arXiv:1706.01749, revised 2019, Section 1.1, displayed symmetric conjecture (2). \url{https://arxiv.org/abs/1706.01749}.
\end{description}
\label{end:statement}
\end{document}
